\section{A \texorpdfstring{$5/4$}{5/4} Lower-Bound Family via Single-Cycle Subdivision}\label{sec:single-cycle}

We consider the following family of instances. Let $n\ge 7$ be odd, let $v_1,\dots,v_n$ be the vertices of $K_n$ (indices modulo $n$), and fix the Hamiltonian cycle
\[
C:=v_1v_2\cdots v_nv_1.
\]
For each $i\in[n]$, subdivide the edge $v_iv_{i+1}$ by a new vertex $x_i$, and let $d$ be the shortest-path metric of the resulting graph. We take $r:=v_1$ as depot. The client set is
\[
V:=\{v_2,\dots,v_n,x_1,\dots,x_n\}.
\]
Every original client $v_2,\dots,v_n$ is active with probability $1$, and each subdividing vertex $x_i$ is active independently with probability $1/2$.

We first record the metric structure. For original vertices,
\[
d(v_i,v_{i+1})=2,
\qquad
 d(v_i,v_j)=1\quad\text{if } j\not\equiv i\pm 1 \pmod n.
\]
For subdividing vertices,
\[
d(x_i,v_i)=d(x_i,v_{i+1})=1,
\qquad
 d(x_i,v_j)=2\quad\text{if } j\notin\{i,i+1\}.
\]
For two subdividing vertices,
\[
d(x_i,x_j)=2
\quad\Longleftrightarrow\quad
x_i\text{ and }x_j\text{ correspond to adjacent edges of }C,
\]
and otherwise $d(x_i,x_j)=3$.

The main result of this section is the following.

\begin{theorem}\label{thm:single-cycle-main}
For every $\varepsilon>0$, there exists a stochastic TSP instance such that both
\[
\frac{\optprior}{\optpost}\ge \frac54-\varepsilon
\qquad\text{and}\qquad
\frac{\optprior}{\optadapt}\ge \frac54-\varepsilon.
\]
More precisely, for every odd $n\ge 7$, the single-cycle subdivision instance above satisfies
\[
\optprior=\frac{15}{8}n,
\qquad
\optpost\le \frac32 n+3,
\qquad
\optadapt\le \frac32 n+18.
\]
Consequently,
\[
\frac{\optprior}{\optpost}
\ge
\frac{(15/8)n}{(3/2)n+3}
\xrightarrow[n\to\infty]{} \frac54
\qquad\text{and}\qquad
\frac{\optprior}{\optadapt}
\ge
\frac{(15/8)n}{(3/2)n+18}
\xrightarrow[n\to\infty]{} \frac54.
\]
\end{theorem}

The proof is split into three parts: first the exact a priori value (Proposition~\ref{prop:single-prior}), then the a posteriori upper bound (Proposition~\ref{prop:single-post}), and finally the adaptive upper bound (Proposition~\ref{prop:single-adapt}).

\subsection*{The a priori value}

\begin{proposition}\label{prop:single-prior}
For every odd $n\ge 7$, the above instance satisfies
\[
\optprior=\frac{15}{8}n.
\]
\end{proposition}

\begin{proof}
Let $T$ be an arbitrary master tour. Since the original vertices are always active, deleting all subdividing vertices from $T$ yields a cyclic ordering
\[
(r=u_1,u_2,\dots,u_n,r)
\]
of the original vertices. For each $j\in[n]$, let
\[
Y_j=(y_{j,1},\dots,y_{j,k_j})
\]
be the ordered sequence of subdividing vertices that appears between $u_j$ and $u_{j+1}$ in $T$, where $u_{n+1}:=r$. Then
\[
\sum_{j=1}^n k_j=n.
\]

For an ordered pair of original vertices $u,v$ and an ordered sequence $Y=(y_1,\dots,y_k)$ of distinct subdividing vertices, let
\[
\Phi(u;Y;v)
\]
denote the expected length of the realized path from $u$ to $v$ obtained from the segment
\[
(u,y_1,\dots,y_k,v)
\]
after deleting inactive subdividing vertices. By linearity of expectation,
\[
\E[c(T|_A)]
=
\sum_{j=1}^n \Phi(u_j;Y_j;u_{j+1}).
\]
So it suffices to lower-bound each segment.

\medskip
\noindent\textbf{Step 1: a convenient exact formula.}
Set
\[
z_0:=u,\qquad z_i:=y_i\ (1\le i\le k),\qquad z_{k+1}:=v.
\]
For $0\le i<j\le k+1$, let $E_{ij}$ be the event that $z_i$ and $z_j$ are retained and all vertices strictly between them are deleted. Then $z_i$ and $z_j$ become consecutive after shortcutting if and only if $E_{ij}$ occurs, hence
\[
\Phi(u;Y;v)=\sum_{0\le i<j\le k+1} d(z_i,z_j)\Pr(E_{ij}).
\]
Since each subdividing vertex is retained independently with probability $1/2$,
\[
\Pr(E_{0,k+1})=2^{-k},\qquad
\Pr(E_{0,t})=2^{-t},\qquad
\Pr(E_{t,k+1})=2^{-(k-t+1)},\qquad
\Pr(E_{ij})=2^{-(j-i+1)}.
\]
Therefore
\begin{align*}
\Phi(u;Y;v)
&=
2^{-k}d(u,v)
+\sum_{t=1}^k 2^{-t}d(u,y_t)
+\sum_{t=1}^k 2^{-(k-t+1)}d(y_t,v) \\
&\qquad
+\sum_{1\le i<j\le k}2^{-(j-i+1)}d(y_i,y_j).
\end{align*}

Introduce the indicators
\[
\delta:=\mathbf 1_{\{u,v\text{ consecutive on }C\}},
\qquad
\alpha_t:=\mathbf 1_{\{y_t\text{ is incident with }u\}},
\qquad
\beta_t:=\mathbf 1_{\{y_t\text{ is incident with }v\}},
\]
and, for $1\le i<j\le k$,
\[
\gamma_{ij}:=\mathbf 1_{\{y_i,y_j\text{ correspond to adjacent edges of }C\}}.
\]
Then
\[
d(u,v)=1+\delta,
\qquad
 d(u,y_t)=2-\alpha_t,
\qquad
 d(y_t,v)=2-\beta_t,
\qquad
 d(y_i,y_j)=3-\gamma_{ij}.
\]
Substituting and simplifying the geometric sums gives
\[
\Phi(u;Y;v)=1+\frac{3k}{2}+2^{-k}\delta-D(Y),
\]
where
\[
D(Y):=
\sum_{t=1}^k 2^{-t}\alpha_t
+\sum_{t=1}^k 2^{-(k-t+1)}\beta_t
+\sum_{1\le i<j\le k}2^{-(j-i+1)}\gamma_{ij}.
\]

\medskip
\noindent\textbf{Step 2: lower bounds for short segments.}
We claim that every segment with $k:=|Y|$ satisfies
\[
\Phi(u;Y;v)\ge
\begin{cases}
1, & k=0,\\[1mm]
2, & k=1,\\[1mm]
\dfrac{11}{4}, & k=2,\\[2mm]
\dfrac{29}{8}, & k=3.
\end{cases}
\]
Moreover, equality for $k=2$ is attained by $(v_i,x_i,x_{i+1},v_{i+1})$, and equality for $k=3$ is attained by $(v_{i+1},x_{i+1},x_i,x_{i-1},v_i)$.

For $k=0$, we simply have $\Phi(u;\emptyset;v)=d(u,v)\ge 1$.

For $k=1$, say $Y=(y)$, the exact formula gives
\[
\Phi(u;y;v)=\frac12 d(u,v)+\frac12\bigl(d(u,y)+d(y,v)\bigr).
\]
If $u$ and $v$ are consecutive on $C$, then $d(u,v)=2$ and $d(u,y)+d(y,v)\ge 2$. If they are not consecutive, then $d(u,v)=1$ and no subdividing vertex is incident with both endpoints, so $d(u,y)+d(y,v)\ge 3$. In either case, $\Phi(u;y;v)\ge 2$.

For $k=2$, the formula becomes
\[
\Phi(u;y_1,y_2;v)=4+\frac\delta4-D(Y).
\]
If $\delta=0$, then each $y_t$ is incident with at most one of $u,v$, so $\alpha_t+\beta_t\le 1$ for $t=1,2$. Hence
\[
D(Y)
=\frac12\alpha_1+\frac14\alpha_2+\frac14\beta_1+\frac12\beta_2+\frac14\gamma_{12}
\le 1+\frac14=\frac54,
\]
which yields $\Phi(u;y_1,y_2;v)\ge 11/4$. If $\delta=1$, then at most one of $y_1,y_2$ can be the subdivider of the edge $uv$; therefore the endpoint contribution is at most $5/4$, and again $\gamma_{12}\le 1$, so $D(Y)\le 3/2$. Hence
\[
\Phi(u;y_1,y_2;v)\ge 4+\frac14-\frac32=\frac{11}{4}.
\]
Equality is attained by $(v_i,x_i,x_{i+1},v_{i+1})$.

For $k=3$, the formula becomes
\[
\Phi(u;y_1,y_2,y_3;v)=\frac{11}{2}+\frac\delta8-D(Y),
\]
where
\begin{align*}
D(Y)
&=\frac12\alpha_1+\frac14\alpha_2+\frac18\alpha_3
+\frac18\beta_1+\frac14\beta_2+\frac12\beta_3 \\
&\qquad +\frac14\gamma_{12}+\frac18\gamma_{13}+\frac14\gamma_{23}.
\end{align*}
If $\delta=0$, then each $y_t$ is incident with at most one of $u,v$, so the endpoint part is at most
\[
\frac12+\frac14+\frac12=\frac54,
\]
and among three cycle edges there are at most two adjacent pairs, whence
\[
\frac14\gamma_{12}+\frac18\gamma_{13}+\frac14\gamma_{23}\le \frac12.
\]
Thus $D(Y)\le 7/4$, and therefore
\[
\Phi(u;y_1,y_2,y_3;v)\ge \frac{11}{2}-\frac74=\frac{15}{4}>\frac{29}{8}.
\]
If $\delta=1$, then at most one of $y_1,y_2,y_3$ can be incident with both endpoints. Hence the endpoint part is at most $3/2$, while the adjacency part is still at most $1/2$. Thus $D(Y)\le 2$, and so
\[
\Phi(u;y_1,y_2,y_3;v)\ge \frac{11}{2}+\frac18-2=\frac{29}{8}.
\]
Equality is attained by $(v_{i+1},x_{i+1},x_i,x_{i-1},v_i)$.

\medskip
\noindent\textbf{Step 3: long segments.}
Now suppose $k\ge 4$. Since an original vertex is incident with exactly two subdividing vertices,
\[
\sum_{t=1}^k 2^{-t}\alpha_t\le 2^{-1}+2^{-2}=\frac34,
\qquad
\sum_{t=1}^k 2^{-(k-t+1)}\beta_t\le \frac34.
\]
Moreover, the graph on $\{y_1,\dots,y_k\}$ in which two vertices are adjacent when the corresponding cycle edges are adjacent has maximum degree at most $2$, hence at most $k$ edges. Since every coefficient $2^{-(j-i+1)}\le 1/4$,
\[
\sum_{1\le i<j\le k}2^{-(j-i+1)}\gamma_{ij}\le \frac{k}{4}.
\]
Therefore
\[
D(Y)\le \frac34+\frac34+\frac{k}{4}=\frac32+\frac{k}{4},
\]
and so
\[
\Phi(u;Y;v)
\ge
1+\frac{3k}{2}-\left(\frac32+\frac{k}{4}\right)
=
\frac54k-\frac12.
\]
Since $k\ge 4$,
\[
\frac54k-\frac12\ge 1+\frac{7k}{8}.
\]
Hence every segment with $k\ge 1$ satisfies
\[
\Phi(u;Y;v)\ge 1+\frac{7k}{8},
\]
while the case $k=0$ gives $\Phi(u;\emptyset;v)\ge 1$.

\medskip
\noindent\textbf{Step 4: summing over segments.}
Using $\sum_j k_j=n$,
\begin{align*}
\E[c(T|_A)]
&=
\sum_{j=1}^n \Phi(u_j;Y_j;u_{j+1}) \\
&\ge
\sum_{j:k_j=0} 1 + \sum_{j:k_j\ge 1}\left(1+\frac{7k_j}{8}\right) \\
&=
 n + \frac78\sum_{j=1}^n k_j \\
&=
 n+\frac78n
=
\frac{15}{8}n.
\end{align*}
Since $T$ was arbitrary, this proves $\optprior\ge 15n/8$.

\medskip
\noindent\textbf{Step 5: matching upper bound for odd $n$.}
Write $n=2m+1$ with $m\ge 3$. Consider the master tour
\begin{align*}
T^\star:=(&v_n,v_2,x_2,x_1,x_n,v_1,
  v_4,x_4,x_3,v_3,
  v_6,x_6,x_5,v_5,
  \dots, \\
 &v_{2m},x_{2m},x_{2m-1},v_{2m-1},v_n).
\end{align*}
This tour contains exactly one $3$-segment,
\[
(v_2,x_2,x_1,x_n,v_1),
\]
whose expected contribution is $29/8$, and exactly $m-1$ $2$-segments,
\[
(v_{2t},x_{2t},x_{2t-1},v_{2t-1})
\qquad (2\le t\le m),
\]
whose expected contributions are all $11/4$. The remaining $m+1$ segments are direct edges between original vertices, each of length $1$. Therefore
\begin{align*}
\E[c(T^\star|_A)]
&=
\frac{29}{8}+(m-1)\cdot \frac{11}{4}+(m+1) \\
&=
\frac{30m+15}{8}
=
\frac{15(2m+1)}{8}
=
\frac{15}{8}n.
\end{align*}
Hence $\optprior\le 15n/8$, and the proof is complete.
\end{proof}

\subsection*{The a posteriori upper bound}

For a realization $A$, let
\[
S(A):=\{i\in[n]: x_i\text{ is active in }A\}
\]
and write $m:=|S(A)|$.

\begin{proposition}\label{prop:single-post}
For every realization $A$ of the single-cycle subdivision instance,
\[
\opt(A)\le n+|S(A)|+3.
\]
Consequently,
\[
\optpost\le \frac32 n+3.
\]
\end{proposition}

\begin{proof}
Fix a realization $A$, and set $S:=S(A)$ and $m:=|S|$.

If $m=n$, then every subdivider is active, and the full subdivided cycle
\[
v_1,x_1,v_2,x_2,\dots,v_n,x_n,v_1
\]
is a feasible tour of cost $2n=n+m\le n+m+3$. So we may assume $m<n$.

Contract every active edge $v_iv_{i+1}$ of the base cycle $C$; equivalently, contract every maximal active run on $C$ into a single block. The resulting quotient is again a cycle, now on
\[
q:=n-m
\]
blocks. Each block is either a single original vertex or a path of the form
\[
v_a,x_a,v_{a+1},x_{a+1},\dots,x_b,v_{b+1}
\]
coming from a maximal consecutive active run. Traversing such a block from one end to the other costs exactly $2\ell$, where $\ell$ is the number of active subdividers in the block, and summing over all blocks gives total internal cost exactly $2m$.

Now consider only the movement \emph{between} blocks. Label the $q$ blocks cyclically by $w_1,\dots,w_q$. If two blocks are consecutive on this quotient cycle, then the distance between the corresponding boundary vertices is at most $2$. If they are not consecutive, then the corresponding boundary originals are nonconsecutive on the original cycle $C$, so their distance is exactly $1$. Thus the inter-block metric is dominated by the standard $q$-vertex cycle metric in which consecutive pairs have distance $2$ and all other pairs have distance $1$.

Let $\tau(q)$ denote the optimal TSP value in that standard $q$-vertex cycle metric. Then the discussion above shows that
\[
\opt(A)\le 2m+\tau(q).
\]
We now bound $\tau(q)$. For $q\ge 5$, there is a Hamiltonian cycle using only nonconsecutive pairs, hence $\tau(q)=q$. Indeed:
\begin{itemize}
    \item if $q$ is odd, the step-$2$ cycle
    \[
    1,3,5,\dots,q,2,4,\dots,q-1,1
    \]
    has total cost $q$;
    \item if $q$ is even, the cycle
    \[
    1,3,5,\dots,q-1,2,q,q-2,\dots,4,1
    \]
    also has total cost $q$.
\end{itemize}
For the small cases, one checks directly that
\[
\tau(1)\le 2,\qquad \tau(2)=4,\qquad \tau(3)=6,\qquad \tau(4)=6.
\]
In particular,
\[
\tau(q)\le q+3 \qquad \text{for every } q\ge 1.
\]
Therefore
\[
\opt(A)\le 2m+(q+3)=2m+(n-m)+3=n+m+3.
\]
Since $m=|S(A)|$, this proves the first claim.

Taking expectation and using $\E[|S(A)|]=n/2$, we obtain
\[
\optpost
=\E_A[\opt(A)]
\le
\E_A[n+|S(A)|+3]
=
 n+\frac n2+3
=
\frac32 n+3.
\]
\end{proof}

\subsection*{The adaptive upper bound}

\begin{proposition}\label{prop:single-adapt}
For every odd $n\ge 7$, the above single-cycle subdivision instance satisfies
\[
\optadapt\le \frac32 n+18.
\]
\end{proposition}

\begin{proof}
We define an explicit adaptive policy $\pi$ and bound its expected cost.

Let
\[
M:=\{v_3,\dots,v_{n-1},x_3,\dots,x_{n-2}\}
\]
be the set of \emph{middle} clients, and let
\[
B:=V\setminus M=\{v_2,v_n,x_1,x_2,x_{n-1},x_n\}
\]
be the set of \emph{boundary} clients.

\medskip
\noindent\textbf{Step 1: visit all boundary clients first.}
The policy first calls all six vertices of $B$ in an arbitrary fixed order, and then returns to the depot $r$.

Every distance in this metric is at most $3$, so this boundary phase has cost at most
\[
7\cdot 3=21,
\]
independently of the realization.

\medskip
\noindent\textbf{Step 2: the alternating middle policy.}
After finishing the boundary phase, the policy is back at $r$. It then calls $v_3$, moves to $v_3$, and from that point on works only on the middle path
\[
P:=v_3,x_3,v_4,\dots,x_{n-2},v_{n-1}.
\]

At every stage, the uncalled middle clients form a subpath of $P$. The current position is one of the two endpoint originals of this remaining subpath. If the current position is the left endpoint original, the policy probes rightward: it calls the adjacent subdividing vertex. If that subdividing vertex is active, then the policy immediately calls the next original vertex and continues in the same direction. If that subdividing vertex is inactive, then the policy jumps to the right endpoint original of the remaining subpath and proceeds symmetrically leftward. The same rule is applied when the current position is the right endpoint original.

Equivalently, the policy keeps moving along the current side of the remaining subpath for as long as the encountered middle subdividing vertices are active, and each time it meets an inactive middle subdividing vertex it switches to the opposite side of the remaining subpath.

\medskip
\noindent\textbf{Step 3: deterministic cost of the middle phase.}
Fix a realization $A$, and let
\[
a_M:=|A\cap\{x_3,\dots,x_{n-2}\}|
\]
be the number of active middle subdividing vertices. Since there are exactly $n-4$ middle subdividing vertices, the number of inactive middle subdividing vertices is
\[
(n-4)-a_M.
\]

We claim that the total cost of the middle phase, including the initial move from $r$ to $v_3$ and the final return to $r$, is at most
\[
n+a_M-1.
\]
Indeed:

\begin{itemize}
    \item The initial move from $r=v_1$ to $v_3$ costs $1$.
    \item Each active middle subdividing vertex contributes exactly $2$: one unit to move from an endpoint original to that subdividing vertex, and one unit to move from there to the next original. Hence all active middle subdividing vertices together contribute exactly
    \[
    2a_M.
    \]
    \item Each inactive middle subdividing vertex causes one switch from one end of the remaining middle subpath to the other. Such a switch costs $1$, except that the very last switch may cost $2$ when only one middle original vertex remains uncalled. Therefore all switches together cost at most
    \[
    ((n-4)-a_M)+1.
    \]
    \item After all middle clients have been called, the salesperson is at some original vertex, and returning from any original vertex to $r$ costs at most $1$.
\end{itemize}

Summing these contributions gives
\[
1+2a_M+\bigl((n-4)-a_M+1\bigr)+1=n+a_M-1.
\]
This proves the claim.

\medskip
\noindent\textbf{Step 4: expectation.}
For every realization $A$, the total cost of $\pi$ is at most
\[
21+(n+a_M-1)=n+a_M+20.
\]
Taking expectation and using $\E[a_M]=(n-4)/2$, we obtain
\[
\E[c(T(A;\pi))]
\le n+\frac{n-4}{2}+20
=\frac32 n+18.
\]
Hence
\[
\optadapt\le \frac32 n+18.
\]
\end{proof}

\subsection*{Proof of Theorem~\ref{thm:single-cycle-main}}

Proposition~\ref{prop:single-prior} gives $\optprior=\frac{15}{8}n$. Proposition~\ref{prop:single-post} gives $\optpost\le \frac32 n+3$. Proposition~\ref{prop:single-adapt} gives $\optadapt\le \frac32 n+18$. Hence
\[
\frac{\optprior}{\optpost}
\ge
\frac{(15/8)n}{(3/2)n+3}
\xrightarrow[n\to\infty]{} \frac54
\qquad\text{and}\qquad
\frac{\optprior}{\optadapt}
\ge
\frac{(15/8)n}{(3/2)n+18}
\xrightarrow[n\to\infty]{} \frac54.
\]
Given any $\varepsilon>0$, choosing odd $n$ large enough yields both ratios at least $5/4-\varepsilon$. \qed
